\chapter{辐射源项分析}

\section{源项参数}

\begin{table}[H]
\centering
\begin{tabular}{|c|c|}
\hline
项目 & 参数 \\
\hline
核素 & $^{192}\text{Ir}$ \\
\hline
活度 & $3.7×10^{5}\text{MBq}(10\text{Ci})$ \\
\hline
半衰期 & 74d \\
\hline
γ 平均能量 & 0.37MeV \\
\hline
剂量常数 & $0.111\mu\text{Sv·m}^2/(\text{MBq·h})$ \\
\hline
源分类 & Ⅲ 类密封放射源 \\
\hline
\end{tabular}
\end{table}


\section{运行工况辐射}

\begin{enumerate}
\item \textbf{正常工况}：有用 γ 射线、贮源器泄漏射线、人体 / 墙体散射射线；
\item \textbf{异常 / 事故工况}：卡源无法回源、源脱落遗落患者体内、联锁失效误照射、人员误入机房、源丢失被盗、屏蔽破损漏射线。
\end{enumerate}

\section{职业病危害因素}

\begin{itemize}
\item \textbf{放射性危害}：$^{192}\text{Ir}$产生 γ 外照射；
\item \textbf{非放射危害}：空气电离产生臭氧、氮氧化物。
\end{itemize}

\section{小结}

主要危害为 γ 射线及臭氧、氮氧化物，需屏蔽 + 机械通风管控。
